% !TEX root = ../main.tex

\begin{abstract}[zh]
	火焰熄火是制约航空发动机与燃气轮机燃烧室稳定运行的关键瓶颈。上游来流速度脉动作为燃烧系统中普遍存在的非定常扰动形式，通过改变横向射流的穿透深度、混合特性及火焰结构，直接影响火焰的稳定燃烧状态，是诱发火焰熄火的重要因素之一。然而，目前关于速度脉动对横向射流火焰熄火边界定量影响规律及其内在机理的研究仍不充分，尤其在双喷嘴构型下的相关研究更为有限。
	
	本文以双喷嘴横向射流火焰为研究对象，以上游速度脉动为扰动输入，以声激励为主动调控手段，系统探究上游速度脉动对横向射流火焰熄火边界的影响机理。研究搭建了集成声激励系统的两级串联横向射流燃烧实验平台，设计加工了四组不同孔径（2mm/3mm）与间距（2d/3d）的双喷嘴，并建立了基于双色示踪粒子图像测速（PIV）的流场-混合物分数同步测量方法，实现了燃料流与空气流速度场的独立求解与混合物分数的定量表征，核心技术已申请发明专利。
	
	主要工作与阶段性结果包括：（1）完成了声激励脉动系统的设计、加工与调试，实现了45 Hz-2.5 kHz频率范围内正弦波与方波激励的稳定输出；（2）筛选出一级燃烧室5组稳定工况，出口温度覆盖1200-1500 K，为二级横向射流实验提供了可靠的热来流条件；（3）完成了双色PIV光路系统搭建与数据处理流程建立，采用365nm紫外LED阵列配合正交柱面镜组形成片光，实现了红蓝粒子的自适应分离、速度场求解与混合物分数计算；（4）通过前置实验明确了6mm孔径喷嘴存在燃烧不充分的问题，确立了采用2mm/3mm小孔径双喷嘴进行正式实验的技术路线；（5）冷态流场DMD分析表明，流场能量集中于低阶大尺度模态，为后续声激励下流场主导模态的识别提供了对比基准。
	
	本研究揭示了声激励参数（频率、波形）与上游速度脉动的定量关联，阐明了双喷嘴构型参数对火焰熄火边界的影响规律，并从流场-混合协同诊断的角度建立了“速度脉动→涡结构演变→混合物分数变化→火焰响应”的物理分析框架。研究成果为燃烧室主动控制与熄火边界预测提供了实验依据与理论支撑。
	
	\vspace{1em}
\end{abstract}

\begin{abstract}[en]
	Flame blowout constitutes a critical bottleneck constraining the stable operation of combustors in aviation engines and gas turbines. As a ubiquitous unsteady disturbance in combustion systems, upstream flow velocity fluctuation directly influences the flame stabilization state by altering the penetration depth, mixing characteristics, and flame structure of transverse jets, thereby serving as a key factor in triggering flame blowout. However, the quantitative effects of velocity fluctuation on the blowout boundary of transverse-jet flames and the underlying mechanisms remain insufficiently investigated, particularly in dual-nozzle configurations.
	
	This study focuses on dual-nozzle transverse-jet flames, taking upstream velocity fluctuation as the disturbance input and acoustic excitation as the active control means, to systematically explore the influence mechanisms of upstream velocity fluctuation on the blowout boundary. An experimental platform of a two-stage series transverse-jet combustor integrated with an acoustic excitation system was established. Four sets of dual nozzles with different orifice diameters (2 mm/3 mm) and spacing (2d/3d) were designed and fabricated. A simultaneous flow-field and mixture-fraction measurement method based on two-color tracer particle image velocimetry (PIV) was developed, enabling independent velocity-field solutions for the fuel stream and the air stream as well as quantitative characterization of the mixture fraction; the core technology has been patented.
	
	The main work and interim results include: (1) Design, fabrication, and commissioning of the acoustic excitation system, achieving stable output of sinusoidal and square-wave excitations over a frequency range of 45 Hz–2.5 kHz; (2) Screening of five stable operating conditions for the primary combustor, with exit temperatures covering 1200–1500 K, providing reliable hot crossflow conditions for the secondary transverse-jet experiments; (3) Completion of the two-color PIV optical system setup and data processing workflow, employing a 365-nm ultraviolet LED array combined with orthogonal cylindrical lens groups to form a light sheet, and achieving adaptive separation of red and blue particles, velocity-field solution, and mixture-fraction computation; (4) Preliminary experiments confirming the issue of incomplete combustion for the 6-mm-diameter nozzle, leading to the adoption of 2-mm and 3-mm small-diameter dual nozzles for the formal experimental route; (5) Dynamic mode decomposition (DMD) analysis of the cold flow field revealing that the flow energy is concentrated in low-order large-scale modes, providing a baseline for identifying dominant flow-field modes under acoustic excitation.
	
	This study reveals the quantitative correlations between acoustic excitation parameters (frequency, waveform) and upstream velocity fluctuation, elucidates the effects of dual-nozzle geometric parameters on the flame blowout boundary, and establishes a physical analysis framework of ``velocity fluctuation $\rightarrow$ vortex-structure evolution $\rightarrow$ mixture-fraction variation $\rightarrow$ flame response'' from the perspective of coupled flow-field and mixing diagnostics. The findings provide experimental evidence and theoretical support for active combustor control and blowout-boundary prediction.
	
	\vspace{1em}
\end{abstract}